\documentclass[tikz,border=6pt]{standalone}
\usepackage{ctex}
\usepackage{amsmath}
\usetikzlibrary{calc, arrows.meta}

\begin{document}
\begin{tikzpicture}[scale=1.1, thick]

% 抛物线 y^2 = 2px，取 p=2，即 y^2 = 4x
% 焦点 F(p/2, 0) = (1, 0)，准线 x = -p/2 = -1

% 坐标轴
\draw[->, thin, gray] (-2.2, 0) -- (5.0, 0) node[right, font=\small] {$x$};
\draw[->, thin, gray] (0, -3.2) -- (0, 3.2) node[above, font=\small] {$y$};
\node[font=\footnotesize, below right] at (0,0) {$O$};

% 抛物线：y^2 = 4x，参数化 t ∈ [-2.8, 2.8]，x=t^2/4, y=t
\draw[black, thick, domain=-2.8:2.8, smooth, samples=80]
  plot ({(\x)^2/4}, \x);

% 准线 x = -1
\draw[red, thick] (-1, -3.0) -- (-1, 3.0);
\node[red, font=\small, left] at (-1, 2.6) {准线 $l$};
\node[red, font=\small, below left] at (-1, 0) {$x=-\dfrac{p}{2}$};

% 焦点 F
\fill[black] (1, 0) circle [radius=2.5pt];
\node[below, font=\small] at (1, -0.05) {$F\!\left(\dfrac{p}{2},0\right)$};

% 动点 P（取参数 t=2，得 P=(1, 2)）
\coordinate (P) at (1, 2);
\fill[black] (P) circle [radius=2.5pt];
\node[above right, font=\small] at (P) {$P(x_0, y_0)$};

% A 是 P 在准线上的射影（垂足）
\coordinate (A) at (-1, 2);
\fill[black] (A) circle [radius=2pt];
\node[left, font=\small] at (A) {$A$};

% PF（蓝色虚线）
\draw[blue, dashed, thick] (P) -- (1, 0)
  node[midway, right, font=\footnotesize] {$|PF|$};

% PA（红色，到准线距离）
\draw[red, thick] (P) -- (A)
  node[midway, above, font=\footnotesize] {$|PA|$};

% 直角标记（PA 与准线垂直）
\draw[thin] (-1+0.12, 2) -- (-1+0.12, 2-0.12) -- (-1, 2-0.12);

% 等距标注
\node[font=\small, align=center] at (2.0, -3.7)
  {$|PF| = |PA|$（到焦点 $=$ 到准线）};

% 顶点标注
\fill[black] (0,0) circle [radius=2pt];
\node[below left, font=\footnotesize] at (0.05, 0) {顶点};

\end{tikzpicture}
\end{document}
